% ABSTRAK page — Juknis p.3: heading, title repeated, "Oleh /
% Nama / NIM / Program Studi", then body (500-800 words, Times
% New Roman 12, not bold), keywords (max. 7) separated by two
% blank lines from the body. Word count not independently verified
% in this pass (no word-count tool run against the final prose).
\chapter*{Abstrak}
\addcontentsline{toc}{chapter}{Abstrak}

% Metadata from config.tex — edit that file, not this one.
\begin{center}
{\bfseries\fontsize{14}{14}\selectfont \ThesisTitleID}

\vspace{1\baselineskip}
Oleh\\
\textbf{\ThesisAuthor}\\
NIM: \ThesisNIM\\
(\ThesisProgramID)
\end{center}

\vspace{1\baselineskip}

\TODO{TERJEMAHKAN SETELAH VERSI BAHASA INGGRIS DIFINALISASI. Draf di bawah ini adalah terjemahan kerja, bukan final.}

Sistem Pengenalan Ekspresi Wajah (\textit{Facial Expression Recognition}, FER) umumnya dikembangkan dan dievaluasi pada data acuan yang kurang merepresentasikan karakteristik wajah dan warna kulit tertentu, sehingga menimbulkan pertanyaan mengenai kemampuan generalisasi sistem tersebut pada populasi di luar distribusi data latihnya. Tesis ini menginvestigasi pertanyaan tersebut menggunakan sebuah dataset wajah masyarakat Papua yang diekstraksi dari sebuah film dokumenter, melalui tiga pertanyaan penelitian (RQ1--RQ3; lihat Bab 1). Evaluasi model FER terlatih (HSEmotion) pada 227 potongan wajah yang dilabel secara manual menemukan akurasi keseluruhan sebesar 22,47\% (RQ1); pada 85 dari 227 potongan wajah yang memiliki data warna kulit (proksi luminansi, bukan ukuran demografis tervalidasi), ditemukan selisih akurasi sebesar 29,71 poin persentase antara dua subkelompok warna kulit yang dikarakterisasi, disertai kecenderungan kuat salah mengklasifikasikan ekspresi Netral sebagai Marah. Tahap praproses CLAHE adaptif yang diuji ditemukan menurunkan, bukan meningkatkan, akurasi klasifikasi pada sampel 227 ini dalam satu kali uji coba eksploratif tanpa uji signifikansi statistik (RQ2). Perbandingan retrospektif pada subset evaluasi bersama berjumlah 135 sampel (lima kelas), antara klasifier HSEmotion dan klasifier Regresi Logistik yang dilatih di atas \textit{embedding} wajah ArcFace yang dibekukan, tidak menemukan perbedaan akurasi keseluruhan yang signifikan secara statistik (uji McNemar, $p = 0{,}46$), meskipun pola kesalahan kedua model berbeda secara kualitatif (RQ3). Secara paralel, tesis ini mendokumentasikan upaya yang masih berlangsung dan bersifat pendukung untuk membangun perluasan dataset beranotasi titik-fitur wajah (Papuan V2); upaya anotasi ini baru mencapai tahap usulan otomatis dan belum ditinjau oleh manusia pada saat penulisan, serta tidak digunakan sebagai bukti bagi RQ1--RQ3.

\vspace{2\baselineskip}
Kata kunci: pengenalan ekspresi wajah, HSEmotion, ArcFace, CLAHE, warna kulit (proksi).
